\documentclass[../readings.tex]{subfiles}

\begin{document}

\subsection{Linearity}
\label{sub:linearity}

\textBox{%
Linear systems are the basic computational problem in scientific computing. Even when the original problem is nonlinear, algorithms often reduce it to a sequence of linear systems by linearizing the problem near the current estimate.
}

\textBox{%
A matrix equation $\bA\bx=\bb$ asks whether the data vector $\bb$ can be produced by combining the columns of $\bA$. The unknown vector $\bx$ contains the weights in that combination. Existence, uniqueness, and numerical difficulty are all controlled by the geometry of those columns.
}

\subsubsection{Linear Systems}

\textBox{%
A linear system is a collection of linear equations involving the same set of variables. In matrix-vector notation, we package these equations into a single compact form that is amenable to both theoretical analysis and high-performance computation.
}

\addDef{Linear System}{%
The matrix-vector form of a linear system is \textbf{$\bA\bx = \bb$}, where:
\begin{itemize}
  \item $\bA \in \fR^{m \times n}$ is the coefficient matrix representing the system's structure.
  \item $\bx \in \fR^n$ is the vector of unknown variables we wish to determine.
  \item $\bb \in \fR^m$ is the right-hand side vector, often representing observed data or a source term.
\end{itemize}
}
\label{def:linear-system}

\addRemark{%
\textbf{(Shape discipline)}\enspace If $\bA \in \fR^{m\times n}$, then $\bx$ must have $n$ entries and $\bb$ must have $m$ entries. The number of columns is the number of unknowns. The number of rows is the number of equations or measurements.
}

\addExmpl{%
\textbf{(Two equations)}\enspace The system
\begin{align*}
    x + 2y = 3, \qquad 3x-y = 1
\end{align*}
can be written as
\begin{align*}
    \begin{pmatrix}1&2\\3&-1\end{pmatrix}
    \begin{pmatrix}x\\y\end{pmatrix}
    =
    \begin{pmatrix}3\\1\end{pmatrix}.
\end{align*}
The first row encodes the first equation; the second row encodes the second equation.
}

\addExcr{%
\begin{enumerate}
    \item Convert $\{x + 2y = 3,\; 3x - y = 1\}$ to the form in Def~\ref{def:linear-system}.
    \item Check the dimensions of $\bA$, $\bx$, and $\bb$ before solving. Which part of Def~\ref{def:linear-system} determines each shape?
    \item Solve in NumPy via \texttt{np.linalg.solve(A, b)}. Verify with the residual from Def~\ref{def:residual}.
\end{enumerate}
}
\label{ex:linear-system-first-solve}

\subsubsection{Rows, Columns, and Geometry}

\textBox{%
There are two complementary ways to read $\bA\bx=\bb$.
\begin{itemize}
    \item \textbf{Row view:}\enspace Each row is one linear equation. In two or three dimensions, each equation describes a line, plane, or hyperplane. Solving the system means finding their intersection.
    \item \textbf{Column view:}\enspace The product $\bA\bx$ is a weighted sum of the columns of $\bA$. Solving the system means expressing $\bb$ as a linear combination of those columns.
\end{itemize}
}

\addRemark{%
\textbf{(The Column Perspective)}\enspace
If $\bA=[\ba_1,\ba_2,\dots,\ba_n]$, then
\begin{align*}
    \bA\bx = x_1 \ba_1 + x_2 \ba_2 + \dots + x_n \ba_n.
\end{align*}
The system $\bA\bx = \bb$ is consistent exactly when $\bb$ lies in the span of the columns of $\bA$.
}
\label{rem:column-perspective}

\addThm{Superposition Principle}{%
If $\bA\bx_1 = \bb_1$ and $\bA\bx_2 = \bb_2$, then for any scalars $\alpha, \beta$:
\begin{align*}
    \bA(\alpha\bx_1 + \beta\bx_2) = \alpha\bb_1 + \beta\bb_2.
\end{align*}
\textbf{Implication:}\enspace Solutions to $\bA\bx = \bb$ can be decomposed and recombined linearly.
}
\label{thm:superposition}

\addExcr{%
\begin{enumerate}
    \item Use Thm~\ref{thm:superposition} to solve $\bA\bx = 3\bb_1 - 2\bb_2$ given only the solutions $\bx_1, \bx_2$ for $\bb_1, \bb_2$.
    \item How does \texttt{np.linalg.solve(A, B)} for a matrix $\bB$ exploit this?
\end{enumerate}
}
\label{ex:superposition-multiple-rhs}

\subsubsection{Invertibility and Existence}

\textBox{%
For square systems ($\bA \in \fR^{n \times n}$), the most critical question is whether a unique solution exists for any given right-hand side. This property is known as invertibility or nonsingularity.
}

\textBox{%
For rectangular systems, the same questions still matter but the answer changes.
\begin{itemize}
    \item \textbf{Square:}\enspace $m=n$. A unique solution may exist.
    \item \textbf{Overdetermined:}\enspace $m>n$. There are more equations than unknowns; usually no exact solution exists, so we solve a least squares problem.
    \item \textbf{Underdetermined:}\enspace $m<n$. There are fewer equations than unknowns; if a solution exists, it is usually not unique.
\end{itemize}
}

\addThm{Equivalences for Invertibility}{%
The following are equivalent for $\bA \in \fR^{n \times n}$:
\begin{enumerate}
    \item $\bA$ is \textbf{invertible} ($\det(\bA) \neq 0$).
    \item $\bA\bx = \bb$ has a \textbf{unique solution} for every $\bb$.
    \item $\bA\bx = \bzero$ has only the \textbf{trivial solution} ($\bx = \bzero$).
    \item $\text{rank}(\bA) = n$ (Full Rank).
    \item Columns of $\bA$ are \textbf{linearly independent}.
\end{enumerate}
}
\label{thm:invertibility-equivalences}

\addRemark{%
\textbf{(Determinants are not algorithms)}\enspace The determinant is useful theoretically, but numerical codes do not test invertibility by computing $\det(\bA)$. Rank, singular values, and condition numbers are more informative and more stable.
}

\addExcr{%
\begin{enumerate}
    \item Determine invertibility using Thm~\ref{thm:invertibility-equivalences}: $\bA = \begin{pmatrix} 2 & -1 \\ 4 & -2 \end{pmatrix}$, $\bB = \begin{pmatrix} 1 & 0 & 0 \\ 0 & 2 & 0 \\ 0 & 0 & 3 \end{pmatrix}$.
    \item Check consistency of $\bA = \begin{pmatrix}1&2\\3&6\end{pmatrix}, \bb = \begin{pmatrix}1\\4\end{pmatrix}$ using the column perspective in Remark~\ref{rem:column-perspective} and \texttt{np.linalg.matrix\_rank}.
\end{enumerate}
}
\label{ex:invertibility-consistency}

\subsubsection{Residuals and Approximate Solutions}

\addDef{Residual}{%
For an approximate solution $\hat{\bx}$ to $\bA\bx=\bb$, the residual is
\begin{align*}
    \br = \bb-\bA\hat{\bx}.
\end{align*}
It measures how well $\hat{\bx}$ satisfies the equations.
}
\label{def:residual}

\addRemark{%
\textbf{(Residual vs.\ error)}\enspace The error is $\hat{\bx}-\bx^*$, where $\bx^*$ is the exact solution. The residual is computable; the error usually is not. A small residual means the equations are nearly satisfied, but it does not always mean the solution is accurate. Ill-conditioned matrices can turn small residuals into large errors.
}
\label{rem:residual-vs-error}

\addExmpl{%
\textbf{(Inconsistent system)}\enspace The equations
\begin{align*}
    x+y=1, \qquad x+y=3
\end{align*}
cannot both be true. In matrix form, $\bb$ is not in the column space of $\bA$. The best we can do is choose $\hat{\bx}$ so that the residual $\bb-\bA\hat{\bx}$ is small in a chosen norm. This is the least squares problem.
}

\subsubsection{Linear Combinations and Bases}

\addDef{Linear Independence}{%
Vectors $\{\bx_1, \dots, \bx_k\}$ are independent if $\sum \alpha_i \bx_i = \bzero \Rightarrow \alpha_i = 0$ for all $i$.
}
\label{def:linear-independence}
\addDef{Basis}{%
A linearly independent set that spans the space. The number of vectors in any basis is the \textbf{dimension}.
}
\label{def:basis}

\subsubsection{Rank}

\textBox{%
The rank of a matrix is the number of independent directions its columns span. It measures how many genuinely different constraints or directions are present after removing redundancy.
}

\addDef{Rank}{%
The rank, $\text{rank}(\bA)$, is the dimension of the column space. It is always true that $\text{rank}(\bA) \leq \min(m, n)$.
\begin{itemize}
    \item \textbf{Full Rank:} The matrix has the maximum possible rank, $\text{rank}(\bA) = \min(m, n)$.
    \item \textbf{Rank Deficient (Singular):} A square matrix with $\text{rank}(\bA) < n$. Such systems do not have unique solutions.
\end{itemize}
}
\label{def:rank}

\addRemark{%
\textbf{(Rank and solution structure)}\enspace For $\bA\in\fR^{m\times n}$:
\begin{itemize}
    \item Full column rank ($\text{rank}(\bA)=n$) means the columns are independent.
    \item Full row rank ($\text{rank}(\bA)=m$) means the rows are independent.
    \item Rank deficiency means there is redundancy: some row or column information can be produced from the rest.
\end{itemize}
}

\addRemark{%
\textbf{(Numerical Rank)}\enspace In code, do not use exact row-reduction to find rank; it is numerically unstable. Use \texttt{np.linalg.matrix\_rank}, which estimates rank from the singular values. A matrix can be mathematically full rank but numerically rank deficient if one singular value is tiny compared to the others.
}

\addThm{Rank-Nullity}{%
For $\bA \in \fR^{m \times n}$: $\text{rank}(\bA) + \text{nullity}(\bA) = n$.
\begin{itemize}
    \item \textbf{Nullity:} Dimension of the null space $\{\bx : \bA\bx = \bzero\}$.
\end{itemize}
}
\label{thm:rank-nullity}

\addRemark{%
\textbf{(General solution form)}\enspace If $\bA\bx=\bb$ is consistent and $\bx_p$ is one particular solution, then every solution has the form
\begin{align*}
    \bx = \bx_p + \bz, \qquad \bz \in \mathcal{N}(\bA).
\end{align*}
The null space contains the directions you can add without changing $\bA\bx$.
}

\addExcr{%
\begin{enumerate}
    \item Use Thm~\ref{thm:rank-nullity} to explain why a ``wide'' matrix ($m<n$) must have a non-trivial null space.
    \item Find the general solution to a consistent $2 \times 3$ system in the form $\bx=\bx_p+\bz$, where $\bz\in\mathcal{N}(\bA)$.
    \item Check directly that adding a null-space vector leaves $\bA\bx$ unchanged.
\end{enumerate}
}
\label{ex:rank-nullity-general-solution}

\subsubsection{The Four Fundamental Subspaces}

\textBox{%
The structure of any linear operator $\bA \in \fR^{m \times n}$ with rank $r$ is completely characterized by four subspaces. Understanding these spaces is essential for solving least squares problems and understanding the behavior of iterative solvers.
}

\begin{enumerate}
    \item \textbf{Column Space ($\mathcal{R}(\bA)$):} $\fR^m, \dim=r$. Span of columns.
    \item \textbf{Null Space ($\mathcal{N}(\bA)$):} $\fR^n, \dim=n-r$. Solutions to $\bA\bx = \bzero$.
    \item \textbf{Row Space ($\mathcal{R}(\bA^T)$):} $\fR^n, \dim=r$. Span of rows.
    \item \textbf{Left Null Space ($\mathcal{N}(\bA^T)$):} $\fR^m, \dim=m-r$. Solutions to $\bA^T\by = \bzero$.
\end{enumerate}

\addThm{Orthogonal Complements}{%
\begin{itemize}
    \item $\mathcal{R}(\bA^T) \perp \mathcal{N}(\bA)$ in $\fR^n$.
    \item $\mathcal{R}(\bA) \perp \mathcal{N}(\bA^T)$ in $\fR^m$.
\end{itemize}
\textbf{Solvability:}\enspace $\bA\bx = \bb$ has a solution iff $\bb \in \mathcal{R}(\bA)$, i.e., $\bb \perp \mathcal{N}(\bA^T)$.
}
\label{thm:orthogonal-complements}

\addExcr{%
\begin{enumerate}
    \item For $\bA = \begin{pmatrix}1&2&3\\2&4&6\end{pmatrix}$, use Def~\ref{def:rank} and Thm~\ref{thm:rank-nullity} to find dimensions and bases for all four spaces.
    \item Verify the two orthogonality statements in Thm~\ref{thm:orthogonal-complements}.
    \item Verify $\bb = (1, 3)^T$ is not in the column space. What is the projection of $\bb$ onto $\mathcal{R}(\bA)$?
\end{enumerate}
}
\label{ex:four-subspaces-projection}

\end{document}
